\begin{abstract}
\noindent
This is the master reference edition of a four-year program on the
strong-coupling spectral structure of Hamiltonian $\SU(N)$ lattice gauge
theory and on the obligations that separate that structure from a
continuum Yang--Mills mass gap. It is organized so that its two halves can
be read, and used, independently.

Part~\ref{part:exact} is a closed body of exact mathematics. A
rank-one cellular subspace --- the kernel of the plaquette boundary map ---
carries the strong-coupling band, and its hopping amplitude is the
rational function
\[
  t_N = \frac{2N\,(N^2-4)}{(N^2-1)(2N^2-1)(4N^2-9)},
\]
positive for $N \ge 3$, vanishing identically at $N = 2$. The four
shared-link channel weights that produce it are order-two Weingarten
values rather than an isotropy assumption, so the second-order theory
rests on representation theory alone. At fourth order the effective
operator splits into Hodge terms and one geometric term coupling the
carrier to its complement, with the off-axis coefficient decided by
derivation across three independent implementations. The rank grading is
exact: every even order of the one-plaquette splitting carries the same
single power of $N$, with coefficients
$c_2 = -4$, $c_4 = -32$, $c_6 = -1748/3$, $c_8 = -123332/9$,
$c_{10} = -49593808/135$, and a four-channel cancellation that sums to
zero at every even order from four to ten; while the odd orders are
determinant families --- for even $N$ every odd order vanishes
identically, and for odd $N$ the first surviving odd order is $N-2$, with
a closed-form single vertex. At sixth order, a rational band shape is
proved to survive every local direct term. Of the statements in
Part~\ref{part:exact}, none is conditional on an open estimate.

Part~\ref{part:fixed} states the construction at fixed lattice spacing:
an actual infinite-volume Wilson transfer limit, a complete physical odd
band isomorphic to $\ell^2(\Z^3) \otimes \C^3$, and a literal
plaquette-source frame that is onto that entire band, on an explicit
interval of the coupling.

Part~\ref{part:continuum} is the part written for whoever works on this
next. It states six obstructions between the fixed-spacing construction
and the continuum limit --- five of them with explicit finite
counterexamples, four of them having killed arguments this program had
written down as complete --- of which the decisive one is that the
interval where the construction is proved and the trajectory along which
the continuum is taken are disjoint sets of couplings. It then gives the
refutation register: \ContraFalsified{} falsified claims with what
falsified each, and the mechanisms this program advanced and withdrew. It
closes with an obligation register of eight numbered problems and the
repository's full route record --- \RouteTotal{} routes, of which
\RouteDone{} are complete, \RouteLive{} live, \RouteUntried{} untried and
\RouteDead{} dead, each dead one carrying the reason it died.

No claim is made here about the existence of continuum Yang--Mills theory
or about a continuum mass gap. \S\ref{sec:breaks} is the argument that no
such claim currently follows, and \S\ref{sec:obligations} is the list of
what would have to change.

Every count quoted above is generated from the repository's ledgers
rather than typed, and every statement carries its
machine-verification tier and the command that re-establishes it.
Currently: \LeanTotal{} Lean~4 theorems with \LeanSorry{} occurrences of
\texttt{sorry}, \ResTotal{} registered analytic results, and
\LitVerified{} verified evidence edges to published work across
\LitObtained{} obtained papers.
\end{abstract}
